Training creates reusable models; add encodes rows and builds lists; search uses the completed index. A cache hit skips fitting, not database encoding. Cold train-plus-add and cached-state search therefore measure different work.

\begin{figure}[tbp]
\centering
% Audited lifecycle; no proposed storage or distributed subsystem.
\begin{tikzpicture}[
 font=\small\linespread{1}\selectfont,
 box/.style={draw=ThesisNavy!65,rounded corners=2pt,align=center,
   text width=3.65cm,minimum height=1.9cm,inner sep=5pt},
 gpu/.style={box,fill=ThesisNavy!7},
 arr/.style={-{Latex[length=2mm]},draw=ThesisNavy!80,thick}]
\node[box] (host) at (0,0) {Host input vectors\\Training samples\\Database batches};
\node[box] (model) at (4.8,0) {Host model\\QR rotation $\Pi$\\Primary codebook};
\node[box] (cache) at (9.6,0) {Host model cache\\$\Pi$ and codebooks\\IVF centroids};
\node[gpu] (train) at (0,-3.6) {GPU training\\Rotate samples\\Fit IVF centroids\\Fit residual model};
\node[gpu] (add) at (4.8,-3.6) {GPU bulk build\\1 or 2 passes\\Assign / encode\\Reorder / pack\\Floats $\leq768$ MiB};
\node[gpu] (index) at (9.6,-3.6) {GPU-resident index\\Offsets and IDs\\Primary / residual codes\\Corrections};
\node[box] (query) at (0,-7.2) {Host query batch\\Pinned staging\\Transfer to GPU};
\node[gpu] (search) at (4.8,-7.2) {GPU query pipeline\\Rotate; probe; scan\\Refine; final top-$k$\\Workspace / graph};
\node[box] (result) at (9.6,-7.2) {Host query results\\IDs and scores\\Transfer from GPU};
\draw[arr] (host) -- node[left,font=\footnotesize] {samples} (train);
\draw[arr] (model.south west) -- (train.north east);
\draw[arr] (train) -- node[above,font=\footnotesize] {model} (add);
\draw[arr] (host.west) -- ++(-.5,0) |- (4.8,-5.15)
 node[pos=.8,below,font=\footnotesize] {database batches} -- (add.south);
\draw[arr] (model) -- (add);
\draw[arr] (add) -- (index);
\draw[<->,dashed,draw=ThesisNavy!80,thick]
 (train.north east) -- node[pos=.8,above,sloped,font=\footnotesize,fill=white,inner sep=1pt]
 {save / load model} (cache.south west);
\draw[arr] (query) -- (search);
\draw[arr] (search) -- (result);
\draw[arr] (index.south) -- ++(0,-.6) -| (search.north);
\end{tikzpicture}

\caption[Build and query architecture]{Implemented build and query boundaries. Host work includes
sampling, orthogonal-matrix construction and analytical primary-codebook
generation; the evaluated configuration fits IVF and residual models on
the GPU. The cache stores trained state, whereas the completed database
representation resides on the device. Dashed arrows denote trained-state
reuse. The staging label bounds only raw/rotated float buffers at source
defaults; other pipeline scratch is additional.}
\label{fig:4-architecture}
\end{figure}

The host prepares samples, the square transform and primary codebook. GPU products apply rotation; device training fits IVF centroids and residual codebooks. Add then transfers bounded batches, encodes both levels and writes list-ordered arrays.

Model state holds the transform, codebooks and centroids. Database state adds offsets, IDs, codes and corrections. Queries overwrite workspace but reuse its allocations; a captured graph references those existing arrays.

One successful \texttt{add} creates a fixed snapshot; another add is rejected. Updates require a replacement snapshot. Physical positions stay fixed during search; mutable insert/delete and transactional publication are unimplemented.

Search maps host queries to host IDs and computed scores. \Cref{eq:5-refined-constant} explains the query-only score term; \cref{sec:3-3} defines timing.

